\chapter{Translations and momentum sectors}
\label{chap:translations}

For a translationally invariant gapped Hamiltonian, Hastings
\cite[lines 1603--1654]{Hastings2006Solving} decomposes the ground space into
momentum sectors and translates one filtered term far from another within a
supercell, so that their low-energy commutator is small. When the momentum
denominators grow with the system size, a finite-orbit recurrence for the
translations on the operators of the ground space gives the same conclusion
\cite{gap:hastings05_translation_supercell}. A variational bound for the
translates of one term closes the chapter.

\section{Site permutations}
\label{sec:translations_site_permutations}

\begin{definition}[Site permutations]
    \label{def:hastings_site_permutation}
    \lean{ManyBody.sitePermOp, ManyBody.sitePermOp_toPerm_mul,
        ManyBody.sitePermOp_mem_unitary,
        ManyBody.sitePermOp_mul_mul_conjTranspose}
    \leanok
    For $g\in G$, $P_g$ is the permutation matrix of $x\mapsto g\cdot x$ on the
    configurations. The map $g\mapsto P_g$ is a unitary representation of $G$, and
    $(P_gAP_g^\dagger)(x,y)=A(x\circ g,y\circ g)$ for every operator $A$.
\end{definition}

\begin{lemma}[Support of a translated operator]
    \label{thm:hastings_site_permutation_support}
    \lean{ManyBody.sitePermOp_mul_mul_conjTranspose_mem_supportedOperators}
    \leanok
    \uses{def:hastings_site_permutation}
    If $A$ acts on the sites $S$, then $P_gAP_g^\dagger$ acts on $g\cdot S$.
\end{lemma}
\begin{proof}\leanok
    $P_gAP_g^\dagger$ is $A$ with both configurations precomposed with the
    permutation $v\mapsto g\cdot v$ of the sites, which carries the support $S$ to
    $g\cdot S$.
\end{proof}

\begin{definition}[Invariant operators]
    \label{def:hastings_site_invariant}
    \lean{ManyBody.IsSiteInvariant, ManyBody.isSiteInvariant_of_conj}
    \leanok
    \uses{def:hastings_site_permutation}
    An operator $H$ is invariant under the action of $G$ when $P_gH=HP_g$ for every
    $g\in G$. In particular, $H$ is invariant if $P_gHP_g^\dagger=H$ for every
    $g\in G$.
\end{definition}

\begin{lemma}[Sums of translated terms]
    \label{thm:hastings_site_invariant_sum}
    \lean{ManyBody.isSiteInvariant_sum_of_covariant, ManyBody.isSiteInvariant_sum_conj}
    \leanok
    \uses{def:hastings_site_invariant}
    Let $G$ also act on a finite index set $K$, and let the operators $h_k$ satisfy
    $h_{g\cdot k}=P_gh_kP_g^\dagger$. Then $\sum_kh_k$ is invariant. In particular,
    for one operator $H_0$, the sum $\sum_{g\in G}P_gH_0P_g^\dagger$ of its
    translates is invariant.
\end{lemma}
\begin{proof}\leanok
    \uses{def:hastings_site_permutation}
    Conjugating by $P_g$ sends $h_k$ to $h_{g\cdot k}$, and $k\mapsto g\cdot k$
    permutes $K$. For the translates, take $K=G$ acting on itself by left
    multiplication.
\end{proof}

\begin{lemma}[Filtering commutes with translation]
    \label{thm:hastings_filtered_term_translate}
    \lean{ManyBody.gaussianFilteredTerm_conj,
        ManyBody.gaussianFilteredTerm_sitePermOp_conj}
    \leanok
    \uses{def:hastings_site_invariant,def:hastings_gaussian_filtered_term}
    For a unitary $U$ and Hermitian $H$, filtering $UWU^\dagger$ with respect to
    $UHU^\dagger$ gives $U\widetilde WU^\dagger$. If $H$ is invariant, the filtered
    operator of $P_gWP_g^\dagger$ is therefore $P_g\widetilde WP_g^\dagger$.
\end{lemma}
\begin{proof}\leanok
    \uses{def:hastings_site_permutation}
    $e^{iUHU^\dagger t}=Ue^{iHt}U^\dagger$, so the integrand of the filter is
    conjugated by $U$ at every time, and conjugation by $U$ commutes with the
    integral.
\end{proof}

\begin{lemma}[Graph automorphisms preserve distances]
    \label{thm:hastings_site_action_distance}
    \lean{SimpleGraph.edist_smul_of_adj_iff, SimpleGraph.cycleGraph_adj_smul}
    \leanok
    If every $g\in G$ preserves the adjacency of a graph on $V$, then
    $d(g\cdot x,g\cdot y)=d(x,y)$ for all sites $x$, $y$. The rotations of a ring
    of $n$ sites are such an action.
\end{lemma}
\begin{proof}\leanok
    A walk from $x$ to $y$ is carried to a walk of the same length from $g\cdot x$ to
    $g\cdot y$, so $d(g\cdot x,g\cdot y)\le d(x,y)$; apply this to $g^{-1}$ for the
    reverse inequality.
\end{proof}

\section{Momentum sectors and supercells}
\label{sec:translations_supercells}

\begin{lemma}[Compression of a conjugate]
    \label{thm:hastings_supercell_compression}
    \lean{Submodule.starProjection_mul_conj_mul_starProjection,
        ManyBody.sitePermOp_toPerm_inv,
        ManyBody.groundCompression_sitePermOp_conj}
    \leanok
    \uses{def:hastings_site_permutation}
    Let $Q$ be the orthogonal projection onto a closed subspace $K$ of a Hilbert
    space, and let $U^\dagger$ fix every vector of $K$. Then
    $Q(UWU^\dagger)Q=QWQ$ for every operator $W$. In particular
    $P(P_gWP_g^\dagger)P=PWP$ for every $g\in S$.
\end{lemma}
\begin{proof}\leanok
    $U^\dagger Q=Q$, and taking adjoints $QU=Q$; so
    $QUWU^\dagger Q=QWQ$. For $g\in S$ take $U=P_g$: then
    $U^\dagger=P_{g^{-1}}$ and $g^{-1}\in S$.
\end{proof}

\begin{lemma}[Low-energy commutators across a supercell translation]
    \label{thm:hastings_supercell_commutator}
    \lean{ManyBody.groundCompression_commutator_eq_of_mem_supercell,
        ManyBody.norm_groundCompression_commutator_le_of_mem_supercell}
    \leanok
    \uses{def:hastings_site_permutation,def:hastings_site_invariant,def:hastings_gaussian_filtered_term}
    Let $G$ act on an index set, let the terms satisfy
    $h_{g\cdot k}=P_gh_kP_g^\dagger$, and let $M_k$ be their filtered operators.
    For every $g\in S$ and all indices $i$, $j$,
    \begin{align}
        [PM_iP,PM_jP]=[PM_iP,PM_{g\cdot j}P].
    \end{align}
    Consequently, if $B(i,j)$ bounds $\|[PM_iP,PM_jP]\|$ for all $i$ and $j$,
    then $B(i,g\cdot j)$ bounds it too, for every $g\in S$.
\end{lemma}
\begin{proof}\leanok
    \uses{thm:hastings_supercell_compression,thm:hastings_filtered_term_translate}
    Since $H$ is invariant, $M_{g\cdot j}=P_gM_jP_g^\dagger$, and its
    compression is $PM_jP$ by the previous theorem.
\end{proof}

\begin{lemma}[Decay across a supercell translation]
    \label{thm:hastings_supercell_decay}
    \lean{ManyBody.norm_groundCompression_commutator_le_of_supercell_separation}
    \leanok
    \uses{def:hastings_site_permutation,def:hastings_gaussian_filtered_term,thm:hastings05_blecm}
    Let the sites be the vertices of a finite graph, let each term $h_k$ act on
    a set $X_k$, let the terms satisfy $h_{g\cdot k}=P_gh_kP_g^\dagger$, and let
    every eigenvalue of $H=\sum_kh_k$ equal $E_0$ or be at least $E_0+\Delta E$.
    If $g\in S$ and $X_{g\cdot j}$ is disjoint from the sites within distance $l$
    of $X_i$, then $\|[PM_iP,PM_jP]\|$ is at most the bound of
    Lemma~\ref{thm:hastings05_blecm} for the pair $h_i$, $h_{g\cdot j}$ at
    separation $l$, which decays exponentially in $l$.
\end{lemma}
\begin{proof}\leanok
    \uses{thm:hastings_supercell_compression,thm:hastings05_low_energy_commutator,thm:hastings05_blecm}
    The compression of $M_k$ is that of $h_k$, and the compression of
    $h_{g\cdot j}=P_gh_jP_g^\dagger$ is that of $h_j$ by
    Lemma~\ref{thm:hastings_supercell_compression}. Apply the bound for
    separated operators to $h_i$ and $h_{g\cdot j}$.
\end{proof}

\begin{definition}[Momentum sector]
    \label{def:hastings_momentum_sector}
    \lean{ManyBody.momentumSector}
    \leanok
    \uses{def:hastings_site_permutation}
    For a character $\chi$ of $G$, a homomorphism from $G$ to the multiplicative
    group $\mathbb{C}^\times$, the
    momentum sector of $\chi$ is the set of vectors $v$ in the eigenspace of $H$
    for $E_0$ with $P_gv=\chi(g)v$ for every $g\in G$. For the ground states of
    momentum $\vec k_a$ in \cite[lines 1613--1618]{Hastings2006Solving},
    $\chi(\vec x)=e^{i\vec x\cdot\vec k_a}$.
\end{definition}

\begin{lemma}[The momentum sectors span the ground space]
    \label{thm:hastings_momentum_sectors_span}
    \lean{ManyBody.iSup_momentumSector}
    \leanok
    \uses{def:hastings_momentum_sector,def:hastings_site_invariant}
    Let $G$ be abelian and $H$ invariant. Then the eigenspace of $H$ for $E_0$ is
    the sum of the momentum sectors of the characters of $G$.
\end{lemma}
\begin{proof}\leanok
    Since $H$ is invariant, each $P_g$ maps the eigenspace into itself. A
    permutation has finite order, so $P_g^N=1$ for some $N\ge1$, and since
    $X^N-1$ has no repeated roots the restriction of $P_g$ is diagonalizable.
    These restrictions commute, so they are simultaneously diagonalizable: the
    eigenspace is the sum of the joint eigenspaces on which each $P_g$ acts as
    a scalar $\chi(g)$. On a nonzero joint eigenspace,
    $P_{gh}=P_gP_h$ and $P_1=1$ give $\chi(gh)=\chi(g)\chi(h)$ and $\chi(1)=1$,
    so $\chi$ is a character and the joint eigenspace lies in its momentum
    sector.
\end{proof}

\begin{definition}[The supercell of the momenta]
    \label{def:hastings_supercell_of_momenta}
    \lean{ManyBody.supercell}
    \leanok
    \uses{def:hastings_momentum_sector}
    The supercell of $H$ at $E_0$ is the intersection of the kernels of the
    characters of $G$ whose momentum sector is nonzero, a subgroup of $G$.
\end{definition}

\begin{lemma}[The supercell fixes the ground space]
    \label{thm:hastings_supercell_of_momenta}
    \lean{ManyBody.mem_supercell_iff}
    \leanok
    \uses{def:hastings_supercell_of_momenta,def:hastings_site_invariant}
    Let $G$ be abelian and $H$ invariant. Then $g$ lies in the supercell if and
    only if $P_gv=v$ for every $v$ in the eigenspace of $H$ for $E_0$. So the
    supercell is a subgroup $S$ with the property assumed above, the
    translations by whole supercells of
    \cite[lines 1626--1630]{Hastings2006Solving}.
\end{lemma}
\begin{proof}\leanok
    \uses{thm:hastings_momentum_sectors_span}
    If $g$ lies in the supercell, then $P_g$ fixes every vector of each nonzero
    momentum sector, hence of their sum, which is the eigenspace by
    Lemma~\ref{thm:hastings_momentum_sectors_span}. Conversely, if $P_g$ fixes
    the eigenspace and $v\ne0$ lies in the sector of $\chi$, then
    $\chi(g)v=P_gv=v$, so $\chi(g)=1$.
\end{proof}

\begin{lemma}[Momenta of bounded denominator]
    \label{thm:hastings_supercell_pow}
    \lean{ManyBody.pow_mem_supercell}
    \leanok
    \uses{def:hastings_supercell_of_momenta}
    If $\chi(g)^N=1$ for every $g\in G$ and every character $\chi$ with a nonzero
    momentum sector, then $g^N$ lies in the supercell for every $g\in G$.
\end{lemma}
\begin{proof}\leanok
    $\chi(g^N)=\chi(g)^N=1$ for each such $\chi$.
\end{proof}

\begin{lemma}[Distance and cell coordinates]
    \label{thm:hastings_cell_distance_le}
    \lean{SimpleGraph.natAbs_valMinAbs_sub_le_edist}
    \leanok
    Writing $d(x,y)$ for graph distance,
    $|\pi(y)-\pi(x)|_n\le r\,d(x,y)$.
\end{lemma}
\begin{proof}\leanok
    Along an edge the circular distance to $\pi(x)$ changes by at most $r$,
    so a walk of length $m$ gives displacement at most $rm$. Take a shortest
    walk; for disconnected vertices the graph distance is infinite.
\end{proof}

\begin{lemma}[Multiples of $N$ reach far]
    \label{thm:hastings_multiples_reach_far}
    \lean{ZMod.exists_lt_natAbs_valMinAbs_add}
    \leanok
    For $N\ge1$ and $a,b\in\mathbb{Z}/n$ there is $t\ge0$ such that $b+Nt-a$
    has circular distance more than $\lfloor n/2\rfloor-N$.
\end{lemma}
\begin{proof}\leanok
    The residues $b+Nt-a$ for $t\ge0$ meet every arc of $N$ consecutive
    residues. Choose $t$ with $b+Nt-a$ among the $N$ residues ending at
    $\lfloor n/2\rfloor$.
\end{proof}

\begin{lemma}[A supercell translation carries a site far away]
    \label{thm:hastings_supercell_far}
    \lean{ManyBody.exists_mem_supercell_lt_edist_add,
        ManyBody.exists_mem_supercell_lt_edist,
        ManyBody.exists_mem_supercell_lt_edist_add_ring}
    \leanok
    \uses{thm:hastings_supercell_pow,thm:hastings_cell_distance_le,thm:hastings_multiples_reach_far}
    Let $G$ be abelian, let $H$ be invariant, let $\pi$ and $e$ be as above, and
    let $\chi(g)^N=1$ for every $g$ and every character $\chi$ with a nonzero
    momentum sector, where $N\ge1$. For sites $x$, $y$ some $g$ in the supercell
    satisfies $\lfloor n/2\rfloor<r\,d(x,g\cdot y)+N$, so it carries $y$ to
    graph distance more than $(\lfloor n/2\rfloor-N)/r$ from $x$. For the ring,
    and in the theorem below, $r=1$. In this case, whenever
    $l+N\le\lfloor n/2\rfloor$, the translate can be chosen with
    $d(x,g\cdot y)>l$. With $N$ the least common multiple of the
    denominators of the momenta, this is the supercell
    of \cite[lines 1620--1630]{Hastings2006Solving} and the choice of
    $\vec x'$ in lines 1634--1641.
\end{lemma}
\begin{proof}\leanok
    Take $t$ from Lemma~\ref{thm:hastings_multiples_reach_far} for $a=\pi(x)$ and
    $b=\pi(y)$, and $g=(e^t)^N$, which lies in the supercell by
    Lemma~\ref{thm:hastings_supercell_pow}. Then $\pi(g\cdot y)=\pi(y)+Nt$, and
    Lemma~\ref{thm:hastings_cell_distance_le} bounds the distance.
\end{proof}

\begin{lemma}[Low-energy commutators of a translation-invariant lattice]
    \label{thm:hastings_supercell_commutator_small}
    \lean{ManyBody.norm_groundCompression_commutator_le_of_momenta,
        ManyBody.norm_sitePermOp_mul_mul_conjTranspose}
    \leanok
    \uses{def:hastings_supercell_of_momenta,thm:hastings05_blecm}
    Let the terms $h_k$ be indexed by the sites, with $h_k$ acting on a set
    $X_k\ni k$ of diameter at most $R$ in the lattice graph, let $G$ be an abelian
    group of translations with $h_{g\cdot k}=P_gh_kP_g^\dagger$, and let every
    eigenvalue of $H=\sum_kh_k$ equal $E_0$ or be at least $E_0+\Delta E$. Let
    $\chi(g)^N=1$ for every $g$ and every character $\chi$ with a nonzero
    momentum sector, and let $\pi$, $e$ and $n$ be as above. Then for every $l$
    with $2R+l+N\le\lfloor n/2\rfloor$ the low-energy commutator of the filtered
    terms at $i$ and $j$ obeys the bound of
    Lemma~\ref{thm:hastings05_blecm} at separation $l$, with the norms of $h_i$
    and $h_j$. Taking $l=\lfloor n/2\rfloor-N-2R$ makes it exponentially small in
    $n$. This is the conclusion of \cite[lines 1603--1641]{Hastings2006Solving}
    for the filtered terms $\tilde H_i$.
\end{lemma}
\begin{proof}\leanok
    \uses{thm:hastings_site_invariant_sum,thm:hastings_supercell_of_momenta,thm:hastings_supercell_far,thm:hastings_supercell_decay}
    The sum of covariant terms is invariant. By
    Lemma~\ref{thm:hastings_supercell_far} some $g$ in the supercell carries
    $j$ to distance more than $\lfloor n/2\rfloor-N\ge2R+l$ from $i$, so the
    sites within distance $l$ of $X_i$ miss $X_{g\cdot j}$. The supercell fixes the
    ground space by Lemma~\ref{thm:hastings_supercell_of_momenta}, so
    Lemma~\ref{thm:hastings_supercell_decay} applies, and
    $\|h_{g\cdot j}\|=\|P_gh_jP_g^\dagger\|=\|h_j\|$.
\end{proof}

\begin{lemma}[A closed exponential bound for filtered low-energy commutators]
    \label{thm:hastings_supercell_commutator_exp}
    \lean{ManyBody.norm_groundCompression_commutator_le_of_momenta_exp}
    \leanok
    \uses{thm:hastings_supercell_commutator_small,thm:hastings05_blecm_exp}
    Under the hypotheses of Lemma~\ref{thm:hastings_supercell_commutator_small},
    let $\tilde H_i$ be the Gaussian-filtered term with any nonnegative
    filter parameter. Write $K_g(r+1)^{k}$ for the sphere-growth bound and set
    \begin{align}
        v      & =2v_R|b_0|e^{2R}+1,                       &
        \gamma & =\min\bigl(1/4,\Delta E/(8v)\bigr),\notag   \\
        C      & =2(2+4|b_0|e^{2R})v_R^2K_g
        \frac{2^{k}k!e^{1/2}}{1-e^{-1/2}}+6.\notag
    \end{align}
    Then
    \begin{align}
        \|[P\tilde H_iP,P\tilde H_jP]\|
         & \le C\|h_i\|\,\|h_j\|
        e^{-\gamma(\lfloor n/2\rfloor-N-2R)}.
        \label{eq:hastings_supercell_exp}
    \end{align}
    The subtraction in the exponent is real subtraction, and the estimate
    holds for every system size. This is an explicit bound for the filtered
    terms in the supercell argument of
    \cite[lines 1620--1641]{Hastings2006Solving}.
\end{lemma}
\begin{proof}\leanok
    \uses{thm:hastings_site_invariant_sum,thm:hastings_supercell_of_momenta,thm:hastings_supercell_far,thm:hastings_supercell_commutator,thm:hastings_supercell_commutator_small,thm:hastings05_low_energy_commutator,thm:hastings05_blecm_exp}
    If $\lfloor n/2\rfloor\ge N+2R$, take
    $l=\lfloor n/2\rfloor-N-2R$.
    Lemma~\ref{thm:hastings_supercell_far} gives a supercell translation $g$
    such that
    \begin{align}
        d(i,g\cdot j)                            & >\lfloor n/2\rfloor-N=2R+l,\notag \\
        X_{g\cdot j}\cap\bigcup_{x\in X_i}B_l(x) & =\varnothing.\notag
    \end{align}
    The supercell fixes the ground space, and covariance preserves the norm:
    \begin{align}
        [P\tilde H_iP,P\tilde H_{g\cdot j}P]
                         & =[P\tilde H_iP,P\tilde H_jP],\notag   \\
        \|h_{g\cdot j}\| & =\|P_gh_jP_g^\dagger\|=\|h_j\|.\notag
    \end{align}
    Lemma~\ref{thm:hastings05_blecm_exp} now gives
    (\ref{eq:hastings_supercell_exp}).
    Otherwise the exponential in (\ref{eq:hastings_supercell_exp}) is at least
    one. The commutator is at most $2\|h_i\|\,\|h_j\|$, and $C\ge6$.
    The original nonnegative filter parameter is arbitrary, since
    Lemma~\ref{thm:hastings05_low_energy_commutator} gives
    \begin{align}
        P\tilde H_kP=Ph_kP.\notag
    \end{align}
\end{proof}

\begin{lemma}[The truncated terms up to their truncation errors]
    \label{thm:hastings_supercell_commutator_truncated}
    \lean{ManyBody.norm_groundCompression_commutator_le_of_norm_sub_le,
        ManyBody.norm_groundCompression_commutator_truncated_le_of_momenta}
    \leanok
    \uses{thm:hastings_supercell_commutator_small,thm:hastings05_blecm_exp,thm:hastings_gaussian_filtered_spectral}
    In the setting of Lemma~\ref{thm:hastings_supercell_commutator_small}, let
    $M_k$ be the filter of $h_k$ under a truncated Hamiltonian
    $\sum_{k'\in s_k}h_{k'}$, the term of eq.~\texttt{mid} of
    \cite[lines 285--307]{Hastings2006Solving}, and let
    $\lVert\tilde H_i-M_i\rVert\le\varepsilon_i$ and
    $\lVert\tilde H_j-M_j\rVert\le\varepsilon_j$, as eq.~\texttt{apb} provides.
    Then $\lVert[PM_iP,PM_jP]\rVert$ is at most the bound of
    Lemma~\ref{thm:hastings_supercell_commutator_small} plus
    $2\varepsilon_i\lVert h_j\rVert+2\lVert h_i\rVert\varepsilon_j$. The errors of
    eq.~\texttt{apb} do not decrease with $n$; the exponential smallness in $n$
    for the $M_i$ is Theorem~\ref{thm:hastings_supercell_truncated_small}.
\end{lemma}
\begin{proof}\leanok
    Write $A=P\tilde H_iP$, $A'=PM_iP$, $B=P\tilde H_jP$ and $B'=PM_jP$. Then
    $[A',B']=[A,B]+[A'-A,B']+[A,B'-B]$. Compression by $P$ does not increase
    the norm, $\lVert B'\rVert\le\lVert M_j\rVert\le\lVert h_j\rVert$, and
    $\lVert A\rVert\le\lVert h_i\rVert$, so the last two commutators are at most
    $2\varepsilon_i\lVert h_j\rVert$ and $2\lVert h_i\rVert\varepsilon_j$.
\end{proof}

\begin{theorem}[The source's truncated terms]
    \label{thm:hastings_supercell_truncated_small}
    \lean{ManyBody.gaussianFilteredTerm_truncation_smul,
        ManyBody.gaussianFilteredTerm_truncation_smul,
        ManyBody.norm_groundCompression_truncatedTerm_commutator_le_of_momenta}
    \leanok
    \uses{def:hastings_supercell_of_momenta,thm:hastings05_blecm}
    In the setting of Lemma~\ref{thm:hastings_supercell_commutator_small}, let
    $M_k$ be the filter of $h_k$ under $\sum_{k'\in s_k}h_{k'}$, the term of
    eq.~\texttt{mid} of \cite[lines 285--307]{Hastings2006Solving}, with
    covariant truncations $s_{g\cdot k}=g\cdot s_k$, and let $h_k$ and the terms of
    $s_k$ act on a set $K_k\ni k$ of diameter at most $R'$. Then for every $l$
    with $2R'+l+N\le\lfloor n/2\rfloor$ the low-energy commutator of $M_i$ and
    $M_j$ obeys the bound of Lemma~\ref{thm:hastings05_blecm} at separation
    $l$. Taking $l=\lfloor n/2\rfloor-N-2R'$ makes it exponentially small in $n$,
    which is the conclusion of \cite[lines 1603--1641]{Hastings2006Solving} for
    the source's $M_i$.
\end{theorem}
\begin{proof}\leanok
    \uses{thm:hastings_site_invariant_sum,thm:hastings_supercell_compression,thm:hastings_supercell_of_momenta,thm:hastings_supercell_far,thm:hastings05_blecm}
    Covariance of the terms and of the truncations gives
    $M_{g\cdot k}=P_gM_kP_g^\dagger$, since conjugation by $P_g$ commutes with
    the Gaussian filter. Each $M_k$ acts on $K_k$. Take $g$ in the supercell
    carrying $j$ to distance more than $\lfloor n/2\rfloor-N\ge2R'+l$ from $i$, so
    that the sites within distance $l$ of $K_i$ miss $K_{g\cdot j}$. Then
    $PM_{g\cdot j}P=PM_jP$, eq.~\texttt{blecm} applies to $M_i$ and
    $M_{g\cdot j}$, and $\lVert M_{g\cdot j}\rVert=\lVert M_j\rVert$.
\end{proof}

\section{Translated terms}
\label{sec:translations_terms}

\begin{lemma}[A uniform basis bound for the actual truncated terms]
    \label{thm:hastings05_uniform_truncated_basis}
    \lean{ManyBody.exists_groundSpaceBasis_sum_norm_offDiagonalPartFamily_truncated_le}
    \leanok
    \uses{thm:hastings05_recursive_uniform_basis,thm:hastings05_uniform_ground_basis,thm:hastings05_observable_commutator_decay}
    Let $M_i$ be the actual operators of (mid), obtained by Gaussian
    filtering $h_i$ under the Hamiltonian truncated to its chosen local
    region. Suppose these regions have diameter at most $R_M$ and
    cardinality at most $v_M$, contain the support of $h_i$ and all terms
    retained in its generator, and contain the centre $i$.
    Under the original Hamiltonian gap and local geometric hypotheses,
    including the additional sphere bound
    (\ref{eq:hastings05_polynomial_sphere_scope}),
    put $\gamma$ equal to the original interaction's commutator decay rate
    and $C_M=C(R,v_M,b_0,|K_g|,k_g)e^{\gamma(2R_M+1)}$.
    There is an actual orthonormal ground-space basis satisfying
    \begin{align}
        \sum_i o(M_i)
         & \le(n-1)^2D_n(C_M)\,|b_0||K_g|S_{\gamma/2^n,k_g}.
        \label{eq:hastings05_uniform_truncated_basis}
    \end{align}
    For fixed local data and ground-space dimension the bound is
    independent of the number of sites. No inverse spectral spread,
    conditioning assumption or accumulated truncation error occurs. This
    proves a restricted form of (oibu), not its general metric version
    \cite{gap:hastings05_geometric_growth_scope}.
\end{lemma}
\begin{proof}\leanok
    \uses{thm:hastings05_recursive_uniform_basis,thm:hastings05_observable_commutator_decay}
    Each $M_i$ is Hermitian, has norm at most $\|h_i\|$ and acts on its
    prescribed local region. Apply the low-energy commutator estimate to
    these observables by filtering them again under the original
    Hamiltonian. Its velocity uses the original interaction range and
    norm budget; only the observable support size and separation correction
    use $v_M,R_M$. Normalize the ground-space compressions by $\|h_i\|$,
    construct the recursive basis, and sum its anchor rows as in
    Lemma~\ref{thm:hastings05_uniform_ground_basis}.
\end{proof}

\begin{lemma}[Stability of the summed off-diagonal bound]
    \label{thm:hastings05_basis_perturbation}
    \lean{LowEnergyBasis.offDiagonalPartFamily_sub,
        LowEnergyBasis.norm_offDiagonalPartFamily_le_mul_norm,
        LowEnergyBasis.norm_offDiagonalPartFamily_le_add,
        ManyBody.norm_groundCompression_sub_le,
        ManyBody.sum_norm_offDiagonalPartFamily_groundCompression_le}
    \leanok
    \uses{thm:hastings05_simple_spectrum_basis}
    For an orthonormal family $v$ of length $n>0$, let
    $o_v(A)=\|\operatorname{off}_v(A)\|$. Then
    \begin{align*}
        o_v(A)\le o_v(B)+(n-1)\|A-B\|.
    \end{align*}
    Consequently any two finite operator families $M_i,F_i$ satisfy
    \begin{align*}
        \sum_i o_v(PM_iP)\le\sum_i o_v(PF_iP)
        +(n-1)\sum_i\|M_i-F_i\|.
    \end{align*}
    This transfers the error in (apb) to the summed bound in (oibu)
    of~\cite[lines 300--320 and 1678--1708]{Hastings2006Solving}.
\end{lemma}
\begin{proof}\leanok
    Taking the off-diagonal part is linear and has operator norm at most
    $n-1$, by the matrix-element estimate. Apply the triangle inequality to
    $\operatorname{off}_v(B)+\operatorname{off}_v(A-B)$.
    Compression by an orthogonal projection is contractive on differences;
    sum the resulting inequalities over the family.
\end{proof}

\begin{lemma}[The basis bound with a growing localization radius]
    \label{thm:hastings05_basis_coupled_radius}
    \lean{ManyBody.exists_groundSpaceBasis_sum_norm_offDiagonalPartFamily_truncated_perturbation_le}
    \leanok
    \uses{thm:hastings05_uniform_ground_basis,thm:hastings05_basis_perturbation,thm:hastings_truncated_approximation}
    Under the original Hamiltonian hypotheses of
    Lemma~\ref{thm:hastings05_uniform_ground_basis}, let $B_{\rm orig}$ denote
    its uniform summed bound. There is one orthonormal ground-space basis such
    that, for every $r\ge0$ and every truncation family $s_i$, the actual terms
    $M_i=\operatorname{Filter}_{r,\Delta E}(H_{s_i},h_i)$ satisfy
    \begin{align*}
        \sum_i o(M_i)\le B_{\rm orig}+(n-1)\sum_i
        \|\operatorname{Filter}_{r,\Delta E}(H,h_i)-M_i\|.
    \end{align*}
    In particular a size-dependent localization choice with total error at most
    $\varepsilon$ has bound $B_{\rm orig}+(n-1)\varepsilon$. The total-error
    estimate must be supplied by the quantitative truncation theorem; it is
    not inferred from the fixed-radius bound of
    Lemma~\ref{thm:hastings05_uniform_truncated_basis}. The original-term
    basis still requires polynomial sphere growth with size-independent
    constants; the perturbation argument does not remove this source
    restriction\cite{gap:hastings05_geometric_growth_scope}.
\end{lemma}
\begin{proof}\leanok
    Choose the basis from the original terms, and use
    $P\operatorname{Filter}_{r,\Delta E}(H,h_i)P=Ph_iP$ for every $r\ge0$.
    Apply Lemma~\ref{thm:hastings05_basis_perturbation} to the actual truncated
    terms and these first filters under the full Hamiltonian.
\end{proof}

\section{Finite translation orbits}
\label{sec:translations_finite_orbit}

\begin{definition}[Translation on the ground space]
    \label{def:hastings_ground_translation}
    \lean{ManyBody.groundTranslation,
        ManyBody.groundTranslationConjugation}
    \leanok
    Let $H$ be invariant under unitary site translations $U_g$, and let
    $K=\ker(H-E_0)$ be its ground eigenspace. Write $u_g=U_g|_K$.
    The induced translation on the ground-space operator algebra is
    $T_g(B)=u_gBu_g^{-1}$.
\end{definition}

\begin{lemma}[Compression and translation]
    \label{lem:hastings_ground_translation_compression}
    \lean{ManyBody.sitePermOp_toPerm_one, ManyBody.sitePerm_starProjection_commute,
        ManyBody.groundTranslation_apply_coe,
        ManyBody.groundTranslation_symm_apply,
        ManyBody.groundTranslationConjugation_groundSpaceCompression,
        ManyBody.groundTranslationConjugation_pow_compression_of_covariant,
        ManyBody.groundTranslationConjugation_pow_compression}
    \leanok
    Let $P$ be the orthogonal projection onto $K$, and put
    $C_K(W)=(PWP)|_K$. Then $U_gP=PU_g$, $u_g^{-1}=u_{g^{-1}}$, and
    \[
        T_g(C_K(W))=C_K(U_gWU_g^*).
    \]
    Consequently $T_g^k(C_K(W))=C_K(U_{g^k}WU_{g^k}^*)$ for every $k\ge0$.
    If $M_{g\cdot i}=U_gM_iU_g^*$, this also gives
    $T_g^k(C_K(M_i))=C_K(M_{g^k\cdot i})$.
\end{lemma}
\begin{proof}\leanok
    \uses{def:hastings_ground_translation}
    Commutation with $H$ preserves $K$. The inverse translation also
    preserves $K$, so its image is exactly $K$. An isometry transports
    orthogonal projections onto the corresponding image subspaces;
    hence $U_gP=PU_g$. Restriction gives the inverse and conjugation
    identities, and induction gives the identity for powers.
\end{proof}

\begin{lemma}[A finite family of separated translates]
    \label{lem:hastings_finite_translation_spacing}
    \lean{ZMod.coordinate_pow_smul,
        ZMod.natAbs_valMinAbs_add_mul_ge_of_near,
        SimpleGraph.edist_pow_mul_ge_of_near}
    \leanok
    Let $\pi$ be a cell coordinate modulo $N$, and suppose translation
    $e$ increases $\pi$ by one. Fix integers $L>0$ and $m\ge0$ with
    $2(m+1)L\le N$. If the circular coordinate separation of $i,j$ is
    at most $\lfloor L/2\rfloor$, then the separation of
    $i,e^{kL}\cdot j$ is at least $\lfloor L/2\rfloor$ for every
    $1\le k\le m$. If each graph edge changes the coordinate by at
    most one, their graph distance satisfies the same lower bound.
\end{lemma}
\begin{proof}\leanok
    Translation by $e^r$ adds $r$ to the coordinate, by induction.
    Represent the initial difference by a signed residue of absolute
    value at most $\lfloor L/2\rfloor$. Adding $kL$ places it between
    $\lfloor L/2\rfloor$ and $N-\lfloor L/2\rfloor$, by the assumed
    bound on $mL$. This proves the circular separation, including a
    residue crossing zero. Summing the coordinate changes along a
    graph path gives the graph-distance bound.
\end{proof}

\begin{lemma}[A commutator from finitely many translation iterates]
    \label{thm:hastings_finite_translation_commutator}
    \lean{norm_commutator_le_of_finite_translation_orbit_bounds}
    \leanok
    Let $K$ have complex dimension $n$, and let $T$ be a complex linear
    isometry of its operator space. If $C\ge0$ and
    \[
        \|[A,T^k(B)]\|\le C\qquad(1\le k\le n^2),
    \]
    then $\|[A,B]\|\le(2^{n^2}-1)C$.
\end{lemma}
\begin{proof}\leanok
    Put $m=n^2$ and let $p(z)=\sum_{k=0}^m p_kz^k$ be the
    characteristic polynomial of $T$. A nonzero eigenvector and norm
    preservation imply that every root has modulus one. Thus
    $|p_0|=1$, and the coefficient bound for a monic polynomial with
    roots on the unit circle gives
    $\sum_{k=1}^m|p_k|\le2^m-1$.
    Cayley--Hamilton gives
    \[
        B=-\sum_{k=1}^m p_0^{-1}p_k T^k(B).
    \]
    Apply the linear map $B\mapsto[A,B]$ and the triangle inequality.
    The operator space has dimension $n^2$, as used above. When
    $n=0$, all its operators vanish and the conclusion still holds.
\end{proof}

\begin{lemma}[The integer orbit spacing]
    \label{lem:hastings_translation_floor_decay}
    \lean{translation_half_natDiv_lower_bound,
        exp_neg_half_translation_natDiv_le,
        translation_finiteOrbit_rate_pos}
    \leanok
    For natural numbers $N,m$, put $L=\lfloor N/(2(m+1))\rfloor$.
    Then
    \begin{align}
        \frac{N}{4(m+1)} & \le\frac L2+\frac12, \\
        e^{-\gamma L/2}  & \le e^{\gamma/2}
        e^{-\gamma N/(4(m+1))}\qquad(\gamma\ge0).
        \label{eq:hastings_translation_floor_decay}
    \end{align}
    If $\gamma>0$, the rate $\gamma/(4(m+1))$ is positive.
\end{lemma}
\begin{proof}\leanok
    Euclidean division gives $N=2(m+1)L+r$ with $0\le r<2(m+1)$.
    Divide by $4(m+1)$ and exponentiate the resulting inequality.
\end{proof}

\begin{lemma}[The support offset in the translation estimate]
    \label{lem:hastings_translation_support_offset_decay}
    \lean{translation_exp_floor_decay_le}
    \leanok
    \uses{lem:hastings_translation_floor_decay}
    With $L=\lfloor N/(2(m+1))\rfloor$, $\gamma\ge0$ and $\rho\in\mathbb R$,
    \begin{align}
        e^{-\gamma(L/2-2\rho-1)}
        \le e^{\gamma(2\rho+2)}e^{-\gamma N/(4(m+1))}.
        \label{eq:hastings_translation_support_offset_decay}
    \end{align}
\end{lemma}
\begin{proof}\leanok
    \uses{lem:hastings_translation_floor_decay}
    Factor the left side as $e^{\gamma(2\rho+1)}e^{-\gamma L/2}$.
    The rounding loss in
    (\ref{eq:hastings_translation_floor_decay}) gives the prefactor
    $e^{\gamma(2\rho+3/2)}$, which is at most $e^{\gamma(2\rho+2)}$.
\end{proof}

\begin{lemma}[A common estimate for small and large systems]
    \label{lem:hastings_translation_near_far_decay}
    \lean{one_le_two_pow_sub_one,
        translation_finiteOrbit_decay_of_near_far}
    \leanok
    \uses{lem:hastings_translation_floor_decay,lem:hastings_translation_support_offset_decay}
    Let $m\ge1$, $\gamma>0$, $C\ge0$, $\rho,z\in\mathbb R$, and set
    $L=\lfloor N/(2(m+1))\rfloor$. Suppose $z\le C$ and, whenever
    $L/2>2\rho+1$, suppose also
    $z\le(2^m-1)C e^{-\gamma(L/2-2\rho-1)}$. Then for every $N$,
    \begin{align}
        z\le(2^m-1)C e^{\gamma(2\rho+2)}e^{-\gamma N/(4(m+1))}.
        \label{eq:hastings_translation_near_far_decay}
    \end{align}
    This includes $L=0$ and all other small sizes with the same constants.
\end{lemma}
\begin{proof}\leanok
    \uses{lem:hastings_translation_floor_decay,lem:hastings_translation_support_offset_decay}
    The coefficient satisfies $2^m-1\ge1$. For $L/2>2\rho+1$, apply
    (\ref{eq:hastings_translation_support_offset_decay}) to the far estimate.
    Otherwise (\ref{eq:hastings_translation_floor_decay}) gives
    $N/(4(m+1))\le2\rho+2$, so
    $e^{\gamma(2\rho+2)}e^{-\gamma N/(4(m+1))}\ge1$ and the bound follows
    from $z\le C$.
\end{proof}

\begin{lemma}[Translation of one observable]
    \label{lem:hastings_finite_observable_orbit}
    \lean{ManyBody.norm_groundCompression_commutator_le_of_finite_observable_orbit}
    \leanok
    Let $H$ be translation invariant, let $m\ge\max\{1,(\dim K)^2\}$,
    and suppose $2(m+1)L\le N$. If every translate of $B$ whose centre
    is at graph distance at least $\lfloor L/2\rfloor$ from the centre
    of $W$ satisfies $\|[PWP,PU_gBU_g^*P]\|\le C$, where $C\ge0$,
    then $\|[PWP,PBP]\|\le(2^m-1)C$.
\end{lemma}
\begin{proof}\leanok
    \uses{lem:hastings_finite_translation_spacing,thm:hastings_finite_translation_commutator,lem:hastings_ground_translation_compression}
    For initially separated centres use the identity translation.
    Otherwise translate by $e^{kL}$, for $1\le k\le(\dim K)^2$,
    and apply the finite-dimensional recurrence to $T_{e^L}$.
    Norms before and after restriction to $K$ agree. If $L=0$, the
    original pair already satisfies the hypothesis. Finally
    $1\le2^m-1$ absorbs the cases needing no nontrivial translate.
\end{proof}

\begin{lemma}[Exponential smallness under invariant-Hamiltonian translations]
    \label{thm:hastings_finite_translation_exp}
    \lean{ManyBody.norm_groundCompression_commutator_le_finite_translation_exp}
    \leanok
    Suppose translations preserve the interaction graph and the
    Hamiltonian $H$. Let a cell coordinate modulo $N$ change by at
    most one along every graph edge, and let $e$ advance it by one.
    Assume the local generator bounds of
    Lemma~\ref{thm:hastings05_growth_free_blecm_separation}, and let
    the original spectral gap above $K$ be $\Delta>0$.
    If $\dim K\le n_0$, put
    \[
        m=\max\{1,n_0^2\},\qquad
        v=4v_R|b_0|e^{2R}+1,\qquad
        \gamma=\min\{1/4,\Delta/(8v)\}>0.
    \]
    Let $W,B$ act within graph radius $\rho$ of $i,j$, respectively,
    and let the first support have at most $v_{\rm obs}$ sites. Then
    \begin{align}
        \|[PWP,PBP]\|
         & \le (2^m-1)(4v_{\rm obs}+4)\|W\|\|B\|
        e^{\gamma(2\rho+2)}e^{-\gamma N/[4(m+1)]}.
        \label{eq:hastings_finite_translation_exp}
    \end{align}
    There is no condition on the ground-state momenta or on covariance
    of a chosen decomposition of $H$.
\end{lemma}
\begin{proof}\leanok
    \uses{lem:hastings_finite_observable_orbit,thm:hastings05_growth_free_blecm_separation,lem:hastings_translation_near_far_decay}
    Set $L=\lfloor N/[2(m+1)]\rfloor$. A translation preserves the
    second support's radius and its operator norm. If its centre is
    at distance at least $D$ from $i$, the support separation is at
    least $D-2\rho$. The propagation and original-gap estimate bounds
    the commutator by
    $(4v_{\rm obs}+4)\|W\|\|B\|e^{-\gamma(D-2\rho)}$, with natural
    subtraction when the separation is an integer. Apply the
    preceding lemma at $D=\lfloor L/2\rfloor$. The floor and support
    losses are absorbed by the scalar estimates, yielding
    (\ref{eq:hastings_finite_translation_exp}). For small $N$, use
    $\|[PWP,PBP]\|\le2\|W\|\|B\|$ and the same prefactor.
\end{proof}

\begin{lemma}[The actual truncated Gaussian terms without momentum restrictions]
    \label{cor:hastings_finite_translation_truncated}
    \lean{ManyBody.norm_groundCompression_truncated_commutator_le_finite_translation_exp}
    \leanok
    In the preceding setting let $M_i$ be the actual Gaussian filter
    of $h_i$ under the Hamiltonian retaining a prescribed set of
    interactions near $i$. Suppose the initial term and every retained
    interaction act within graph radius $\rho$ of $i$, in a set of
    at most $v_{\rm obs}$ sites, and $\|h_i\|\le J$ for all $i$.
    Then
    \[
        \|[PM_iP,PM_jP]\|
        \le (2^m-1)(4v_{\rm obs}+4)J^2
        e^{\gamma(2\rho+2)}e^{-\gamma N/[4(m+1)]}.
    \]
    The retained sets need not be covariant. This includes the
    physical ball cutoff of the source and the case of momentum
    denominators growing with $N$.
\end{lemma}
\begin{proof}\leanok
    \uses{thm:hastings_finite_translation_exp,thm:hastings_truncated_support,thm:hastings_gaussian_filtered_spectral}
    The retained Hamiltonian and the initial term have the stated
    support, so their Gaussian evolution and its integral have that
    support. Gaussian filtering contracts the norm, giving
    $\|M_i\|\le J$. Apply the theorem with $W=M_i$, $B=M_j$.
\end{proof}

\section{A variational lower bound}
\label{sec:translations_variational}

\begin{lemma}[Iterated commutators of Hermitian operators]
    \label{thm:hastings_commutator_iter_hermitian}
    \lean{Matrix.isHermitian_smul_commutatorCLM_pow,
        Matrix.conjTranspose_commutatorCLM_pow, Matrix.commutatorCLM_pow_succ_apply}
    \leanok
    For Hermitian $H$ and $A$, a real number $a$ and $k\ge0$, the operator
    $a\,i^k\operatorname{ad}_H^k(A)$ is Hermitian.
\end{lemma}
\begin{proof}\leanok
    By induction, $(\operatorname{ad}_H^kA)^\dagger=(-1)^k\operatorname{ad}_H^k(A)$,
    and $\overline{i^k}(-1)^k=i^k$.
\end{proof}

\begin{lemma}[The variational lower bound for unitary conjugates]
    \label{thm:hastings_variational_lower_bound}
    \lean{Matrix.card_mul_le_of_mulVec_eq_smul_of_conj,
        Matrix.card_mul_sInf_spectrum_le_of_conj,
        Matrix.posSemidef_sub_card_smul_of_conj, Matrix.commutatorCLM_pow_conj,
        Matrix.commutatorCLM_pow_sum, Matrix.commutatorCLM_pow_self_succ}
    \leanok
    \uses{thm:hastings_commutator_iter_hermitian}
    Let $H_0$ be an operator, let $U_1,\dots,U_N$ be unitaries, and let
    $H=\sum_iU_iH_0U_i^\dagger$ with every $U_i$ commuting with $H$. For real
    constants $a_1,\dots,a_m$ let
    \begin{align}
        K=H_0+\sum_{k=1}^ma_k\,i^k\operatorname{ad}_H^k(H_0).
    \end{align}
    If $K-\lambda\ge0$, then $H-N\lambda\ge0$; in particular every eigenvalue
    of $H$, and so its ground energy $E_0$, is at least $N\lambda$. This
    generalizes equation (varnum) of \cite{Hastings2006Solving}, whose terms are
    the translates of one term; the source's form is
    Lemma~\ref{thm:hastings_variational_lower_bound_sites}.
\end{lemma}
\begin{proof}\leanok
    Conjugation by a unitary $U$ that commutes with $H$ commutes with
    $\operatorname{ad}_H$, so
    $\sum_iU_i\operatorname{ad}_H^k(H_0)U_i^\dagger=\operatorname{ad}_H^k(H)=0$
    for $k\ge1$, and $\sum_iU_iKU_i^\dagger=H$. Since $U_i(K-\lambda)U_i^\dagger\ge0$
    and $U_iU_i^\dagger=1$, summing gives $H-N\lambda\ge0$. For an eigenvector
    $\Omega$ of $H$ with eigenvalue $E$,
    $0\le\langle\Omega,(H-N\lambda)\Omega\rangle=(E-N\lambda)\|\Omega\|^2$.
\end{proof}

\begin{lemma}[The variational lower bound]
    \label{thm:hastings_variational_lower_bound_sites}
    \lean{ManyBody.card_mul_le_of_mulVec_eq_smul_of_sitePermOp,
        ManyBody.card_mul_sInf_spectrum_le_of_sitePermOp}
    \leanok
    \uses{thm:hastings_variational_lower_bound,thm:hastings_site_invariant_sum}
    Let a finite group $G$ act on a finite set of sites, such as the
    translations of a periodic lattice of any dimension or a group of
    automorphisms of a finite graph, and let $P_g$ be the unitary that moves the
    configurations of the sites along $g$. For the sum
    $H=\sum_{g\in G}P_gH_0P_g^\dagger$ of the translates of one term $H_0$, and
    $K$ as in Lemma~\ref{thm:hastings_variational_lower_bound}, every
    eigenvalue of $H$ is at least $|G|\lambda$ whenever $K-\lambda\ge0$. For
    Hermitian $H_0$, with $\lambda=\lambda_{\min}(K)$ and $N=|G|$,
    \begin{align}
        E_0/N\ge\lambda_{\min}\Bigl(H_0+\sum_{k=1}^ma_k\,i^k
        \operatorname{ad}_H^k(H_0)\Bigr).
    \end{align}
    Here $N=|G|$ counts the summands. When $G$ acts simply transitively on the
    sites, as the translations of a periodic lattice do, $N$ is the number of
    sites. This is equation (varnum) of \cite{Hastings2006Solving}, lines
    856--877, for translation-invariant $H=\sum_iH_i$; there the terms are
    written with $[H_i,H]=-\operatorname{ad}_H(H_i)$, which only changes the
    signs of the arbitrary constants.
\end{lemma}
\begin{proof}\leanok
    \uses{def:hastings_site_permutation}
    The map $g\mapsto P_g$ is multiplicative, so conjugation by $P_h$ permutes
    the translates and fixes $H$, and each $P_h$ is a unitary commuting with $H$.
    Lemma~\ref{thm:hastings_variational_lower_bound} applies. For Hermitian
    $K$, its real spectrum is finite, and every spectral value is at least
    its infimum. The spectral theorem therefore gives
    $K-\lambda_{\min}(K)\ge0$, yielding the stated specialization.
\end{proof}

\begin{lemma}[The Gaussian-filtered variational lower bound]
    \label{thm:hastings_gaussian_variational_lower_bound}
    \lean{ManyBody.sum_conj_gaussianFilteredTerm,
        ManyBody.card_mul_sInf_spectrum_gaussianFilteredTerm_le_of_conj,
        ManyBody.card_mul_sInf_spectrum_gaussianFilteredTerm_le_of_sitePermOp}
    \leanok
    \uses{thm:hastings_variational_lower_bound,def:hastings_gaussian_filtered_term}
    Let $H_0$ be Hermitian, let $U_1,\ldots,U_N$ be unitaries, and suppose
    each $U_i$ commutes with $H=\sum_i U_iH_0U_i^\dagger$.
    Let $K$ be the Gaussian average of $e^{itH}H_0e^{-itH}$, with variance
    $q/\Delta E^2$, where $q\ge0$ and $\Delta E>0$. Then
    $H=\sum_iU_iKU_i^\dagger$, and every eigenvalue $E$ of $H$ satisfies
    $E\ge N\lambda_{\min}(K)$. For a simply transitive translation action,
    $N$ is the number of sites, so $E_0/N\ge\lambda_{\min}(K)$.
    This is the filtered choice discussed in \cite{Hastings2006Solving},
    source lines 1005--1009, in the variational bound (varnum).
\end{lemma}
\begin{proof}\leanok
    \uses{thm:hastings_variational_lower_bound,thm:hastings_gaussian_filtered_spectral,thm:hastings_filtered_term_translate}
    Covariance of the Gaussian average gives
    $U_iKU_i^\dagger=\widetilde{U_iH_0U_i^\dagger}$, because $U_i$ fixes $H$.
    The filtered terms sum to $H$. Apply
    Lemma~\ref{thm:hastings_variational_lower_bound} with local term $K$
    and all commutator coefficients zero.
\end{proof}
